\documentclass{article}

	\usepackage{graphicx}

	\usepackage{url}

	\usepackage{hyperref}

	\hypersetup{colorlinks=true,urlcolor=blue}

	\usepackage{enumitem}

	\usepackage{geometry}

	\geometry{left =2cm, right=2cm, top=2cm,bottom=2cm}

	\usepackage{xcolor}

	\definecolor{myorange}{RGB}{255, 128, 0}

	\usepackage{titlesec}

	\newcommand{\underlinedtitle}[1]{\color{myorange}#1\par\nobreak\noindent\rule{\textwidth}{1pt}}

	\titleformat{\section}{\normalfont\Large\bfseries}{}{0em}{\underlinedtitle}

	\titlespacing\section{0pt}{12pt plus 4pt minus 2pt}{5pt plus 2pt minus 2pt}

	\titlespacing\subsection{0pt}{12pt plus 4pt minus 2pt}{5pt plus 2pt minus 2pt}

	\titlespacing\subsubsection{0pt}{12pt plus 4pt minus 2pt}{5pt plus 2pt minus 2pt}

	\usepackage{enumitem}

	\setlist[itemize]{noitemsep, topsep=0pt}

	\usepackage{fontawesome}

	\usepackage{orcidlink}

	\renewcommand{\familydefault}{\sfdefault}


	
	\begin{document}
\noindent {\Huge Nelly Barret}

\noindent {\huge Assistant professor in computer science}

\noindent {\Large LIRIS $\cdot$ INSA Lyon}

\noindent \faMapMarker\ Blaise Pascal building, 7 Avenue Jean Capelle, 69100 Villeurbanne, France

\noindent \faBuilding\ Office 502.3.21 $|$ \faEnvelope\ nelly.barret@insa-lyon.fr $|$ \faLaptop\ https://perso.liris.cnrs.fr/nbarret/ $|$ \orcidlink{0000-0002-3469-4149} \href{https://orcid.org/0000-0002-3469-4149}{0000-0002-3469-4149}

\section*{\faHeart\ Research interests}
\noindent My research themes lie in the broad area of heterogeneous data integration and exploitation, including heterogeneous and multi-modal data as well as warehouse, data lake and lakehouse architectures.

\noindent The general scientific questions that are driving me every day include:

	\begin{itemize}

		\item How to effectively and efficiently collect, organize and store heterogeneous data produced by various actors?

		\item How to explore and exploit large amounts of data, especially for domain experts?

		\item How to clean, join, merge, and sementically enrich raw data for better decision making?

	\end{itemize}

	My research applies to various domains including sustainable cities, media, and healthcare, with a strong interest in sustainable cities.\section*{\faBriefcase\ Academic positions}
\noindent\textbf{Assistant Professor} $\cdot$ INSA Lyon \& LIRIS (FR) \hfill \textbf{Sep 2025 - now}
\begin{itemize}
\item Research: part of the DRIM team, focusing on document engineering and distributed systems.
\item Teaching: 1st and 2nd year courses at the engineer school INSA Lyon.
\end{itemize}
\vspace{1em}\noindent\textbf{Postdoctoral researcher} $\cdot$ Politecnico di Milano (IT) \hfill \textbf{Apr 2024 - Jul 2025}
\begin{itemize}
\item Horizon Europe project: "Better real-world health-data distributed analytics research platform".
\item Close collaborators: Pietro Pinoli (associate professor, Politecnico di Milano, IT), Anna Bernasconi (assistant professor, Politecnico di Milano, IT), Boris Bikbov (medical doctor, Politecnico di Milano, IT).
\end{itemize}
\vspace{1em}\noindent\textbf{PhD student} $\cdot$ Institut Polytechnique de Paris \& Inria Saclay (FR) \hfill \textbf{Jan 2021 - Mar 2024}
\begin{itemize}
\item PhD thesis title: "User-oriented exploration of semi-structured datasets".
\item PhD defense jury: Fatiha Sais (présidente – professor, Univ. Paris-Saclay and LISN), Jean-Marc Petit (rapporteur – professor, INSA Lyon and LIRIS), Olivier Teste (rapporteur – professor, Univ. Toulouse Jean Jaurès and IRIT), Katja Hose (examinatrice – professor, TU Wien, AT), Stefano Ceri (examinateur – professor, Politecnico di Milano, IT), Fatemeh Nargesian (examinatrice – associate professor, Univ. Rochester, USA), Ioana Manolescu (advisor – research director, Ecole Polytechnique and Inria Saclay), Karen Bastien (co-advisor – WeDoData CEO).
\end{itemize}
\vspace{1em}\noindent\textbf{Master intern} $\cdot$ LIRIS (FR) \hfill \textbf{Feb 2020 - Jul 2020}
\begin{itemize}
\item Master thesis title (in French): "Prédiction de l'environnement d'un quartier".
\item Advisors: Fabien Duchateau (associate professor, Univ. Lyon 1 and LIRIS) and Franck Favetta (associate professor, Univ. de Lyon and LIRIS).
\item Other collaborators: Nelly Duong (CEO, Home in Love), Behnaz Jullien (intern in psychology, Univ. Lyon 2), Wissame Laddada (post-doctoral researcher, LIRIS), and Ludovic Moncla (associate professor, computer science at INSA Lyon).
\end{itemize}
\vspace{1em}\noindent\textbf{Bachelor intern} $\cdot$ LIRIS (FR) \hfill \textbf{May 2018 - Jul 2018}
\begin{itemize}
\item Bachelor thesis title (in French): "Intégration de données géographiques pour la recommandation de quartiers".
\item Advisors: Fabien Duchateau (associate professor, Univ. Lyon 1 and LIRIS) and Franck Favetta (associate professor, Univ. de Lyon and LIRIS).
\item Other collaborators: Nelly Duong (CEO, Home in Love), Loic Bonneval (associate professor in sociology, Univ. Lyon 2) and Aurélien Gentil (PhD student in sociology, Univ. Lyon 2).
\end{itemize}
\section*{\faGraduationCap\ Education}
\noindent\textbf{PhD degree} $\cdot$ Institut Polytechnique de Paris \hfill \textbf{Jan 2021 - Mar 2024}


\noindent\textit{Track}: Computer Science, Data and Artificial Intelligence

\vspace{1em}\noindent\textbf{Master's degree} $\cdot$ Université Claude Bernard Lyon 1 \hfill \textbf{Sep 2018 - Jul 2020}


\noindent\textit{Track}: Computer Science, Artificial Intelligence

\vspace{1em}\noindent\textbf{Bachelor's degree} $\cdot$ Université Claude Bernard Lyon 1 \hfill \textbf{Sep 2015 - Jul 2018}


\noindent\textit{Track}: Computer Science

\section*{\faLightbulbO \ Research projects}
\noindent\textbf{MOBI4ALL: Towards intelligent systems for inclusive and privacy-respecting mobility for individuals with Autism Spectrum Disorder (ASD)} \hfill \textbf{Oct 2026 - Oct 2030}


\noindent\textit{Role}: member (WP4) $|$ \textit{Grant}: ANR PRC (2026-2030)

\vspace{1em}\noindent\textbf{EXPERT-MED-IA: Le confrère virtuel d'aide à la décision pour le soignant} \hfill \textbf{Oct 2026 - Oct 2029}


\noindent\textit{Role}: member $|$ \textit{Grant}: AURA region (2026-2028)

\vspace{1em}\noindent\textbf{BETTER: Real-world health-data distributed analytics research platform} \hfill \textbf{Apr 2024 - Jul 2025}


\noindent\textit{Role}: contributor (tasks 4.1, 4.2) $|$ \textit{Grant}: Horizon Europe (2022 - 2027) $|$ \textit{Parnters}: 7 European hospitals $|$ \textit{Website}: \href{https://better-health-project.eu}{better-health-project.eu}

\begin{itemize}
\item \textbf{Scientific outcomes:} Two conceptual models for data and metadata respectively, both general, able to represent various multi-modal healthcare data, allowing the usage of ontologies, and extensible to various healthcare scenarios; An ETL algorithm to fully automatically convert existing data and metadata to instances of our models, resulting in an interoperable database; A metadata and data catalogue for exploration and ease federated learning algorithm design.
\item \textbf{Practical outcomes:} Each hospital has an interoperable database, part of the global BETTER network; Federated learning algorithms are co-created and implemented by researchers and practitioners; The catalogue implementation for the 3 use-cases (genetic rare diseases)
\end{itemize}
\vspace{1em}\noindent\textbf{CONNECTIONSTUDIO: User data exploration of heterogeneous data} \hfill \textbf{Mar 2023 - Mar 2024}


\noindent\textit{Role}: contributor $|$ \textit{Grant}: ANR SourcesSay project (2020 - 2024) $|$ \textit{Website}: \href{https://connectionstudio.inria.fr}{connectionstudio.inria.fr}

\begin{itemize}
\item \textbf{Scientific outcomes:} A user-oriented methodology for querying a relational data lake by selecting data pieces of interest instead of formulating SQL queries
\item \textbf{Practical outcomes:} A software suite for the data exploration of the underlying data lake
\end{itemize}
\vspace{1em}\noindent\textbf{PATHWAYS: Find interesting entity-to-entity paths in heterogeneous data} \hfill \textbf{May 2022 - Mar 2024}


\noindent\textit{Role}: contributor $|$ \textit{Grant}: DIM RFSI PHD 2020-01 $|$ \textit{Website}: \href{https://team.inria.fr/cedar/projects/pathways}{team.inria.fr/cedar/projects/pathways}

\begin{itemize}
\item \textbf{Scientific outcomes:} Efficiently enumerate all paths connecting named entities appearing in multi-modal and heterogeneous data by relying on an intermediate summary graph and a view-based algorithm; The ranking of paths based on their interestingness, a quantitative measure based on entity confidence and information dilution
\item \textbf{Practical outcomes:} A novel ChatGPT-based module for entity extraction; The efficient evaluation of paths by using a dedicated multi-query optimization algorithm
\end{itemize}
\vspace{1em}\noindent\textbf{ABSTRA: Entity-Relationship summaries from semi-structured data} \hfill \textbf{Jan 2021 - Mar 2024}


\noindent\textit{Role}: main contributor $|$ \textit{Grant}: DIM RFSI PHD 2020-01 $|$ \textit{Parnters}: WeDoData start-up $|$ \textit{Website}: \href{https://team.inria.fr/cedar/projects/abstra}{team.inria.fr/cedar/projects/abstra}

\begin{itemize}
\item \textbf{Scientific outcomes:} A novel model-agnostic data summarization method relying on data kinds instead of specific model features; The selection of most representative entities, their attributes and relationships from the data summary by relying on a graph representation and PageRank node scores
\item \textbf{Practical outcomes:} A end-to-end pipeline to summarize any (set of) heterogeneous datasets; An interface to visualize summaries as Entity-Relationship diagrams
\item \textbf{Collaboration:} WeDoData is a French SME specializing in data journalism and data visualization.
\end{itemize}
\vspace{1em}\noindent\textbf{PREDIHOOD: Supervised prediction of the environment of neighborhoods} \hfill \textbf{Feb 2020 - Jul 2020}


\noindent\textit{Role}: main contributor $|$ \textit{Grant}: LabEx IMU (Intelligences des Mondes Urbains) $|$ \textit{Parnters}: HomeInLove start-up

\begin{itemize}
\item \textbf{Scientific outcomes:} An algorithm to automatically select the most useful subset of indicators among the hundreds collected ones; An algorithm to automatically predict in a supervised manner the environment of any French neighborhood based on several subsets of indicators
\item \textbf{Practical outcomes:} Collect and process cartographic and urban data from public data such as INSEE; An interface for predicting the environment of any French neighborhood; A general interface for testing, evaluating and comparing AI algorithms
\item \textbf{Collaboration:} HomeInLove is a French start-up helping people in case of professional mobilities to find a new place to live based on their criteria and needs (neighborhood, family, hobbies, ...).
\end{itemize}
\vspace{1em}\noindent\textbf{GEOALIGN: Matching and merging geographic entities} \hfill \textbf{Mar 2019 - Jul 2019}


\noindent\textit{Role}: main contributor $|$ \textit{Grant}: Master 1 project

\begin{itemize}
\item \textbf{Scientific outcomes:} A customizable similarity formula for estimating the matching similarity between geographic entities; The computation of the estimated quality of the matching without ground truth
\item \textbf{Practical outcomes:} A map interface for find, matching and merging similar geographic entities
\end{itemize}
\vspace{1em}\noindent\textbf{VIZLIRIS: Recommending and clustering French neighborhoods} \hfill \textbf{May 2019 - Jun 2019}


\noindent\textit{Role}: main contributor $|$ \textit{Grant}: LabEx IMU (Intelligences des Mondes Urbains) $|$ \textit{Parnters}: HomeInLove start-up $|$ \textit{Website}: \href{https://perso.liris.cnrs.fr/fabien.duchateau/?page=VizLIRIS/}{perso.liris.cnrs.fr/fabien.duchateau/?page=VizLIRIS/}

\vspace{1em}\noindent\textbf{FLAME: Federated LeArning in Multimodal Environments} \hfill \textbf{Jan 1970}


\noindent\textit{Role}: leader $|$ \textit{Grant}: ANR JCJC (2026-2030)

\section*{\faCode\ Research tools}
\noindent\textbf{I-ETL} \hfill \textbf{Apr 2024 - Jul 2025}


\noindent\textit{Role}: developer $|$ \textit{Language}: Python $|$ \textit{LOC}: 6k $|$ \textit{Repository}: \href{https://github.com/DEIB-GECO/i-etl}{github.com/DEIB-GECO/i-etl}

\vspace{1em}\noindent\textbf{PathWays} \hfill \textbf{May 2022 - Mar 2024}


\noindent\textit{Role}: main developer $|$ \textit{Language}: Java $|$ \textit{LOC}: 4k $|$ \textit{Repository}: \href{https://gitlab.inria.fr/cedar/pathways}{gitlab.inria.fr/cedar/pathways}

\vspace{1em}\noindent\textbf{Abstra} \hfill \textbf{Jan 2021 - Mar 2024}


\noindent\textit{Role}: developer $|$ \textit{Language}: Java $|$ \textit{LOC}: 10k $|$ \textit{Repository}: \href{https://gitlab.inria.fr/cedar/abstra}{gitlab.inria.fr/cedar/abstra}

\vspace{1em}\noindent\textbf{ConnectionStudio} \hfill \textbf{Mar 2023 - Mar 2024}


\noindent\textit{Role}: contributor $|$ \textit{Language}: Java and Javascript $|$ \textit{LOC}: 25k $|$ \textit{Repository}: \href{https://gitlab.inria.fr/cedar/connection-studio}{gitlab.inria.fr/cedar/connection-studio}

\vspace{1em}\noindent\textbf{Predihood} \hfill \textbf{Feb 2020 - Jul 2020}


\noindent\textit{Role}: main developer $|$ \textit{Language}: Python and JavaScript $|$ \textit{LOC}: 3k $|$ \textit{Repository}: \href{https://gitlab.com/fduchate/predihood}{gitlab.com/fduchate/predihood}

\vspace{1em}\noindent\textbf{GeoAlign} \hfill \textbf{Jan 2019 - Jun 2019}


\noindent\textit{Role}: main developer $|$ \textit{Language}: PHP and JavaScript $|$ \textit{LOC}: 4k $|$ \textit{Repository}: private

\vspace{1em}\noindent\textbf{VizLIRIS} \hfill \textbf{May 2018 - Jul 2018}


\noindent\textit{Role}: main developer $|$ \textit{Language}: Python and JavaScript $|$ \textit{LOC}: 1k $|$ \textit{Repository}: \href{https://forge.univ-lyon1.fr/stage-Nelly/HiL-recommender}{forge.univ-lyon1.fr/stage-Nelly/HiL-recommender}

\section*{\faBook\ Publications}
\subsection*{International journals}
\begin{enumerate}[resume]
\item Nelly Barret, Sourav Bhowmick, Angela Bonifati, Barbara Catania, Stratos Idreos, Ekaterini Ioannou, Sana Sellami, Madhulika Mohanty, Roee Shraga, Utku Sirin, Juno Steegmans, Pinar Tözün, Soror Sahri, Genoveva Vargas-Solar. \href{https://sigmodrecord.org/publications/sigmodRecord/2509/pdfs/08_Reports_Barret.pdf}{Diversity, Equity and Inclusion Activities in Database Conferences: A 2024 Report}. ACM SIGMOD Record. 2025
\item \underline{Nelly Barret}, Anna Bernasconi, Boris Bikbov, Pietro Pinoli. \href{https://www.researchgate.net/publication/396456781_I-ETL_an_interoperability-aware_health_metadata_pipeline_to_enable_federated_analyses}{I-ETL: an interoperability-aware health (meta)data pipeline to enable federated analyses}. BMC Medical Informatics and Decision Making. 2025
\item \underline{Nelly Barret}, Antoine Gauquier, Jia-Jean Law, Ioana Manolescu. \href{https://hal.science/hal-04727209}{Finding meaningful paths in heterogeneous graphs with PathWays}. Information Systems. 2025
\item \underline{Nelly Barret}, Fabien Duchateau, Franck Favetta, Aurélien Gentil, Loïc Bonneval. \href{https://hal.science/hal-05388570}{An Environmental Study of French Neighbourhood}. Communications in Computer and Information Science. 2021
\item \underline{Nelly Barret}, Fabien Duchateau, Franck Favetta. \href{https://joss.theoj.org/papers/10.21105/joss.02805}{Predihood: an open-source tool for predicting neighbourhoods' information}. Journal of Open Source Software. 2021
\end{enumerate}
\subsection*{International conferences}
\begin{enumerate}[resume]
\item \underline{Nelly Barret}, Anna Bernasconi, Cinzia Cappiello, Giacomo Palu, Pietro Pinoli. \href{https://www.researchgate.net/publication/392478298_Leveraging_Profiling_to_Bridge_Healthcare_Silos_for_Federated_Analyses}{Leveraging profiling to bridge healthcare silos for federated analyses}. Conference on Advanced Information Systems Engineering. 2025
\item \underline{Nelly Barret}, Ioana Manolescu, Prajna Upadhyay. \href{https://openproceedings.org/2024/conf/edbt/paper-31.pdf}{Computing generic abstractions from application datasets}. Extending DataBase Technology conference. 2024
\item \underline{Nelly Barret}, Antoine Gauquier, Jia-Jean Law, Ioana Manolescu. \href{https://hal.science/hal-04131977}{Exploring heterogeneous data graphs through their entity paths}. Advances in DataBases and Information Systems. 2023
\item Nelly Barret, Simon Ebel, Théo Galizzi, Ioana Manolescu, Madhulika Mohanty. \href{https://hal.science/hal-04185938}{User-friendly exploration of highly heterogeneous data lakes}. International Conference on Cooperative Information Systems. 2023
\item \underline{Nelly Barret}, Ioana Manolescu, Prajna Upadhyay. \href{https://dl.acm.org/doi/10.1145/3511808.3557179}{Abstra: toward generic abstractions for data of any model}. International Conference on Information and Knowledge Management. 2022
\item \underline{Nelly Barret}, Fabien Duchateau, Franck Favetta, Loïc Bonneval. \href{https://www.scitepress.org/PublicationsDetail.aspx?ID=YRZoV7ChM6U=&t=1}{Predicting the environment of a neighbourhood: a use case for France}. International Conference on Data Science, Technology and Applications. 2020
\end{enumerate}
\subsection*{International workshops}
\begin{enumerate}[resume]
\item \underline{Nelly Barret}, Mélina Verger. \href{https://hal.science/hal-05579539}{SemiStructQuest: Unveiling NoSQL database concepts through gamified learning platform analyses}. Data Education workshop. 2026
\item Nelly Barret, Tudor Enache, Ioana Manolescu, Madhulika Mohanty. \href{https://inria.hal.science/hal-04591933}{Finding the PG schema of any (semi) structured dataset: a tale of graphs and abstraction}. SEAGRAPH workshop. 2024
\item Oana Balalau, Simon Ebel, Théo Galizzi, Ioana Manolescu, Madhulika Mohanty. \href{https://inria.hal.science/hal-04591897}{Graph lenses over any data: the ConnectionLens experience}. SEAGRAPH workshop. 2024
\end{enumerate}
\subsection*{Demonstrations}
\begin{enumerate}[resume]
\item \underline{Nelly Barret}, Antoine Gauquier, Jia-Jean Law, Ioana Manolescu. \href{https://hal.science/hal-04131977}{PathWays: entity-focused exploration of heterogeneous data graphs}. The Semantic Web: ESWC Satellite Events. 2023
\item \underline{Nelly Barret}. \href{https://hal.science/hal-03344102/document}{Facilitating heterogeneous dataset understanding}. Gestion de Données - Principes, Technologies et Applications. 2021
\item \underline{Nelly Barret}, Fabien Duchateau, Franck Favetta, Ludovic Moncla. \href{https://hal.science/hal-05388560}{Spatial entity matching with GeoAlign}. ACM International Conference on Advances in Geographic Information Systems. 2019
\end{enumerate}
\subsection*{National conferences}
\begin{enumerate}[resume]
\item \underline{Nelly Barret}, Simon Ebel, Théo Galizzi, Ioana Manolescu, Madhulika Mohanty. \href{https://hal.science/hal-04388609}{Exploration utilisateur de lacs de données très hétérogènes}. Conférence Francophone sur l'Extraction et la Gestion des Connaissances. 2024
\item \underline{Nelly Barret}, Fabien Duchateau, Franck Favetta, Maryvonne Miquel, Aurélien Gentil, Loïc Bonneval. \href{https://editions-rnti.fr/?inprocid=1002526}{À la recherche du quartier idéal}. Conférence Francophone sur l'Extraction et la Gestion des Connaissances. 2019
\end{enumerate}
\subsection*{Manuscripts}
\begin{enumerate}[resume]
\item \underline{Nelly Barret}. \href{https://theses.fr/2024IPPAX001}{User-oriented exploration of semi-structured datasets}. Institut Polytechnique de Paris. 2024
\item \underline{Nelly Barret}. \href{https://zenodo.org/records/17679906}{Prédiction de l'environnement d'un quartier}. Université Claude Bernard Lyon 1. 2020
\item \underline{Nelly Barret}. \href{https://zenodo.org/records/17679892}{Alignement d'entités spatiales avec GeoAlign}. Université Claude Bernard Lyon 1. 2019
\item \underline{Nelly Barret}. \href{https://zenodo.org/records/17679854}{Intégration de données géographiques pour la recommandation de quartiers}. Université Claude Bernard Lyon 1. 2018
\end{enumerate}
\section*{\faTrophy\ Awards}
\noindent\textbf{PhD thesis award} $\cdot$ BDA conference \hfill \textbf{Oct 2024}


\noindent\textit{Website}: \href{https://bda2024.sciencesconf.org/resource/page/id/10}{bda2024.sciencesconf.org/resource/page/id/10}

\begin{itemize}
\item My PhD work won the 2nd prize for the PhD award delivered by the BDA conference, awarding one or two PhD with significant contributions in the French data management community.
\end{itemize}
\section*{\faGroup \ Working groups}
\noindent\textbf{LIRIS commission on gender equality and parity} $\cdot$ LIRIS \hfill \textbf{Nov 2025 - now}


\noindent\textit{Role}: member $|$ \textit{Website}: \href{https://vargas-solar.com/aequitas}{vargas-solar.com/aequitas}

\begin{itemize}
\item The aim of this commission is to work towards building a more equitable scientific community by: (a) increasing knowledge about the scientific contributions of the female community by promoting the visibility of women at LIRIS, (b) increase communication in the lab to raise awareness on F/M equality questions, and (c) studying women's participation in the data science and artificial intelligence workforce.
\end{itemize}
\vspace{1em}\noindent\textbf{FAIRification of genomic annotations} \hfill \textbf{Dec 2024 - now}


\noindent\textit{Role}: member $|$ \textit{Website}: \href{https://www.rd-alliance.org/groups/fairification-genomic-annotations-wg/activity/}{www.rd-alliance.org/groups/fairification-genomic-annotations-wg/activity/}

\begin{itemize}
\item We aim at defining novel methods to make genomic annotations and data more interoperable. I co-lead two actions: (1) define a strategy to publicly and permanently store harmonized genomic metadata with Sveinung Gundersen (computer scientist, ELIXIR) and Evan Christensen (PhD in bio-informatics, Univ. d'Utah, CA); and (2) define new models for representing genomic annotations with Sveinung Gundersen and Adam Wright (researcher, Cancer institute of Ontario, USA).
\end{itemize}
\vspace{1em}\noindent\textbf{Diversity, Equity and Inclusion} \hfill \textbf{Apr 2023 - now}


\noindent\textit{Role}: member $|$ \textit{Website}: \href{https://dbdni.github.io/}{dbdni.github.io/}

\begin{itemize}
\item I lead the SCOUT action which aims at facilitating DEI efforts in the database community. My main proposal is a checklist to easily comply with inclusive submission guidelines; it will be integrated into EasyChair and CMT.
\end{itemize}
\section*{\faGroup\ Service}
\subsection*{Leadership}
\subsubsection*{Conference and workshop organizer}
\begin{itemize}
\item $[$N/A$]$ CORIA-TALN (Conférence en Recherche d'Information et Applications - Traitement Automatique des Langues Naturelles): 2027 (co-organizer)
\end{itemize}
\subsubsection*{Conference and workshop chair}
\begin{itemize}
\item $[$A$]$ CMLS@ER (Conceptual Modeling for Life Sciences workshop): 2026 (chair), 2025 (chair)
\item $[$B$]$ EDBT (Extending DataBase Technology conference): 2027 (DEI chair)
\end{itemize}
\subsubsection*{PC responsabilities}
\begin{itemize}
\item $[$A$]$ CIKM (International Conference on Information and Knowledge Management): 2026
\item $[$A$]$ TGD@EDBT (Transforming Graph Data workshop): 2026
\item $[$A$]$ iSAILS@CAiSE (Information Systems and AI for Life Sciences workshop): 2025
\item $[$B$]$ RCIS (Research Challenges in Information Science): 2026 (forum)
\item $[$C$]$ ADBIS (Advances in DataBases and Information Systems): 2026, 2025
\item $[$N/A$]$ BDA (Gestion de Données - Principes, Technologies et Applications): 2026
\item $[$N/A$]$ DBKDA (Databases, Knowledge, and Data Applications): 2025
\end{itemize}
\subsection*{Review responsabilities}
\subsubsection*{Journals}
\begin{itemize}
\item $[$Q1$]$ BMC Bioinfo (BMC Bioinformatics): 2025
\item $[$Q1$]$ GigaScience (GigaScience): 2025
\item $[$Q1$]$ BMC Med. Inf. and Dec. Making (BMC Medical Informatics and Decision Making): 2024
\item $[$Q1$]$ PLoS One (PLoS One): 2026, 2024
\item $[$Q1$]$ Scientific Reports (Scientific Reports): 2024
\item $[$Q2$]$ CMC (Computers, Materials and Continua): 2025
\item $[$Q2$]$ JDIQ (Journal of Data and Information Quality): 2025
\item $[$N/A$]$ JOSS (Journal of Open Source Software): 2025
\end{itemize}
\subsubsection*{Conferences}
\begin{itemize}
\item $[$Q1$]$ VLDB Proceedings (Proceedings of the Very Large DataBases volume): 2027 (review board member)
\item $[$A*$]$ ICDE (International Conference on Data Engineering): 2025
\item $[$A$]$ CAiSE (Conference on Advanced Information Systems Engineering): 2025
\item $[$A$]$ CIDR (Conference on Innovative Data systems Research): 2023
\item $[$B$]$ EDBT (Extending DataBase Technology conference): 2023
\item $[$B$]$ ICWE (International Conference on Web Engineering): 2021
\item $[$N/A$]$ BDA (Gestion de Données - Principes, Technologies et Applications): 2026, 2025 (demo)
\end{itemize}
\section*{\faMicrophone\ Talks and Dissemination}
\subsection*{Seminar}
\begin{itemize}
\item \href{https://hal.science/hal-05623512}{Multimodal data without borders: integration and exploration to the rescue}. Oct 2025, IRIT, Toulouse (FR).
\item \href{https://hal.science/hal-05623544}{Leveraging profiling to bridge healthcare silos}. Jun 2025, FIL, Lyon (FR).
\item \href{https://zenodo.org/records/17642343}{Healthcare warehouses: from data integration to federated analyses}. May 2025, MADICS, Toulouse (FR).
\item \href{https://zenodo.org/records/17679336}{Integrating and exploring heterogeneous datasets}. Apr 2024, POLIMI, Milan (IT).
\item \href{https://zenodo.org/records/17679827}{Heterogeneous datasets: a tale of integration and exploration}. Jan 2024, LIRIS, Lyon (FR).
\item \href{https://zenodo.org/records/17679729}{User-oriented exploration of semi-structured datasets}. Jan 2024, LIB, Dijon (FR).
\item \href{https://zenodo.org/records/17679385}{User-oriented exploration of semi-structured datasets}. Oct 2023, LISN, Paris (FR).
\end{itemize}
\subsection*{Panel}
\begin{itemize}
\item Open debate on AI for/to teaching. Sep 2026, Lyon (FR).
\item \href{https://www.youtube.com/watch?v=xlCy-D-kpvc}{When artificial intelligence inherits our biases}. Oct 2025, Fete de la science, Lyon (FR).
\end{itemize}
\subsection*{Vulgarization}
\begin{itemize}
\item \href{https://zenodo.org/records/17679761}{From data to journalism}. Jan 2024, LIPPS, Palaiseau (FR).
\item \href{https://zenodo.org/records/17680052}{Artificial intelligence: a tool for journalism}. Jun 2023, CFI, Paris (FR).
\item \href{https://zenodo.org/records/17680102}{ConnectionStudio : from data to journalism}. May 2023, DATAJOURNOS, Paris (FR).
\item \href{https://zenodo.org/records/17679370}{Research in data integration: the ConnectionLens project}. Mar 2022, ECOLE POLYTECHNIQUE, Online.
\end{itemize}
\subsection*{Female empowerment}
\begin{itemize}
\item Scientific women. Nov 2025, Lyon (FR).
\item Girls, Mathematics, and Computer Science Day (RJMI). Nov 2025, Lyon (FR).
\item \href{https://www.deib.polimi.it/ita/eventi/dettagli/2987}{Girls in Computer Science and Engineering (Girls@CSE)}. Sep 2024, POLIMI, Milan (IT).
\end{itemize}
\section*{\faUniversity\ Institutional responsabilities}
\noindent\textbf{DRIM team annual seminar} \hfill \textbf{Sep 2026 - now}


\noindent\textit{Role}: Co-organizer

\begin{itemize}
\item Co-organize the DRIM team's annual 2-day seminar, which allows PhDs and young researchers to discuss their work with the team.
\end{itemize}
\vspace{1em}\noindent\textbf{GRID500 platform LIRIS referent} \hfill \textbf{Jan 2026 - now}


\noindent\textit{Role}: Main referent $|$ \textit{Website}: \href{https://www.grid5000.fr/w/Grid5000:Home}{www.grid5000.fr/w/Grid5000:Home}

\begin{itemize}
\item Manage LIRIS account requests on the national platform Grid'5000.
\end{itemize}
\vspace{1em}\noindent\textbf{DRIM webpage maintainer} \hfill \textbf{Sep 2025 - now}


\noindent\textit{Role}: Main maintainer $|$ \textit{Website}: \href{https://liris.cnrs.fr/equipe/drim}{liris.cnrs.fr/equipe/drim}

\begin{itemize}
\item Manage the institutional webpage of the DRIM team at the LIRIS lab.
\end{itemize}
\vspace{1em}\noindent\textbf{Data Science group LinkedIn page maintainer} \hfill \textbf{Apr 2024 - Jul 2025}


\noindent\textit{Role}: Main maintainer $|$ \textit{Website}: \href{https://linkedin.com/company/polimi-data-science-group}{linkedin.com/company/polimi-data-science-group}

\begin{itemize}
\item Manage the LinkedIn page of the Data Science group at DEIB, Politecnico di Milano.
\end{itemize}
\vspace{1em}\noindent\textbf{Inria PhD representative} \hfill \textbf{Jan 2023 - Apr 2024}


\noindent\textit{Role}: Main representative within Ecole Polytechnique

\begin{itemize}
\item Represent Inria Saclay PhD students in Institut Polytechnique de Paris.
\end{itemize}
\vspace{1em}\noindent\textbf{CEDAR team seminar organizer} \hfill \textbf{Jan 2021 - Mar 2024}


\noindent\textit{Role}: main contributor $|$ \textit{Website}: \href{https://team.inria.fr/cedar/category/seminar}{team.inria.fr/cedar/category/seminar}

\begin{itemize}
\item Organize and disseminate the CEDAR team seminars at Inria Saclay.
\end{itemize}
\vspace{1em}\noindent\textbf{Internship application website} \hfill \textbf{Jan 2021 - Dec 2021}


\noindent\textit{Role}: contributor $|$ \textit{Website}: \href{https://stages.dix.polytechnique.fr/home}{stages.dix.polytechnique.fr/home}

\begin{itemize}
\item Co-develop and deploy the website for internships at École Polytechnique, in collaboration with École Polytechnique IT team.
\end{itemize}
\section*{\faSlideshare\ Teaching}
\subsection*{Reccuring courses}
\begin{itemize}
\item ISN1 - Informatique et Société Numérique 1 (37h, INSA Lyon, Engineer 1st year): algorithmic, Python development
\item ISN2 - Informatique et Société Numérique 2 (35h, INSA Lyon, Engineer 1st year): algorithmic, functional decomposition, introduction to networks, recursion, guided hands-on project
\item ISN4 - Informatique et Société Numérique 4 (24h, INSA Lyon, Engineer 2nd year): object-oriented programming, graphical interfaces, code versioning, hands-on project
\item SOL - Systèmes et Outils Logiciels (15h, INSA Lyon, Engineer 1st year): Microsoft Office and LibreOffice suite, basic Linux commands
\end{itemize}
\subsection*{Guest lectures}
\begin{itemize}
\item Introduction to data lakes (1h, INSA Lyon, Engineer 5th year): 45-minutes course on the data lake architectutre
\item Introduction to scientific research (8h, INSA Lyon, Engineer 4th and 5th year): 1-day seminar and workshops on academic research in France
\end{itemize}
\subsection*{Teaching service}
\begin{itemize}
\item Industrial 1st year internship referee (INSA Lyon): 3-5 students each year
\item First year admission interviews (INSA Lyon): recruitment session for high school students who will join INSA as engineers in the fall
\end{itemize}
\subsection*{Previous courses}
\begin{itemize}
\item CSE203 - computer science project (24h, Ecole Polytechnique, Bachelor 2nd year): research-oriented computer science project
\item CSE204 - Machine Learning (24h, Ecole Polytechnique, Bachelor 2nd year): main concepts, paradigms and algorithms in machine learning
\item INF411 - bases de programmation et d'algorithmique (40h, Ecole Polytechnique, Engineer 2nd year): data structures (hash, trees, graphs), algorithmic, and an introduction to complexity
\item INF371 - mécanismes de la programmation orientée objet (40h, Ecole Polytechnique, Engineer 1st year): OOP in Java
\end{itemize}
\section*{\faCommentsO\ Student supervision}
\subsection*{Engineer students}
\begin{itemize}
\item Côme Perin on ``Apprentissage multimodal au service de la santé des conducteurs'' (summer 2026, ENTPE Vaulx-en-Velin)
\item Ahmed Bel Hadj Youssef on ``Building Persistent Episodic and Semantic Memory for LLM Agents via RAG-based Retrieval'' (summer 2026, INSA Lyon)
\item Lou Reina-Kuntziger on ``MemeAltas: a framenet parser for memes'' (summer 2026, INSA Lyon)
\item Crisitan Vaireanu on ``TODO'' (summer 2026, INSA Lyon)
\item Ahmed Bel Hadj Youssef on ``Design and implementation of a data mesh architecture for domain-oriented document classification'' (fall 2025, INSA Lyon)
\end{itemize}
\subsection*{Master students}
\begin{itemize}
\item Inès Boudjeli on ``Débruitage et sélection automatique de signaux cardiaques pour la prédiction médicale'' (spring 2026, Université de Dijon)
\item Antoine Gauquier on ``Path exploration of ConnectionLens graphs'' (summer 2022, IMT Nord-Europe)
\end{itemize}
\subsection*{Bachelor students}
\begin{itemize}
\item Shay Pripstein on ``Target data acquisition in open repositories'' (winter 2023, École Polytechnique)
\item Tudor Enache on ``From simple to property graphs'' (winter 2023, École Polytechnique)
\item Nikola Dobricic on ``Graph queries for abstractions'' (summer 2023, École Polytechnique)
\item Jia-Jean Law on ``Optimization of entity-to-entity data paths'' (winter 2022, Ecole Polytechnique)
\end{itemize}
\section*{\faWpforms\ Training}
\noindent\textbf{ACM Peer Reviewer Training (1 hour)} \hfill \textbf{May 2026}


\noindent\textit{Website}: \href{https://reviewers.acm.org/training-course}{reviewers.acm.org/training-course}


\noindent Peer-review overview; Assessing your suitability to review; Review touchstones; Evaluating the paper; Submitting your review; Artifact review and badging

\vspace{1em}\noindent\textbf{Doctoral supervision (18 hours)} \hfill \textbf{Apr 2026}


\noindent\textit{Website}: \href{https://fedora.insa-lyon.fr/fr/page/seminaire-encadrants}{fedora.insa-lyon.fr/fr/page/seminaire-encadrants}


\noindent The PhD journey as a system; Human and organizational factors for a successful PhD journey; Scientific co-production between the PhD student and the advisors; Understanding inter-culturality within the PhD eco-system; Evolution of the supervisor's role during a PhD and through a career

\vspace{1em}\noindent\textbf{Student evaluation (4 hours)} \hfill \textbf{Apr 2026}

\noindent How to evaluate students within the APC framework (Competency-based approach)?; How to evaluate students in the GenAI era?


\vspace{3em}\noindent\centering\textcolor{gray}{Last update: \today}

\end{document}